\documentclass[11pt,a4paper]{ctexart}
\usepackage[a4paper,margin=1.78cm,headheight=15pt]{geometry}
\usepackage{amsmath,amssymb,mathtools,bm}
\usepackage{booktabs,longtable,tabularx,array}
\usepackage{enumitem}
\usepackage[most]{tcolorbox}
\usepackage{tikz}
\usetikzlibrary{arrows.meta,positioning,fit,calc}
\usepackage{xcolor}
\usepackage{fancyhdr}
\usepackage{hyperref}
\usepackage{microtype}

\newcolumntype{Y}{>{\raggedright\arraybackslash}X}
\newcolumntype{L}[1]{>{\raggedright\arraybackslash}p{#1}}
\setlength{\parindent}{2em}
\setlength{\parskip}{0.34em}
\setlength{\emergencystretch}{3em}
\renewcommand{\arraystretch}{1.22}
\setlength{\tabcolsep}{3pt}
\setlist[itemize]{itemsep=0.18em,topsep=0.35em,leftmargin=2em}
\setlist[enumerate]{itemsep=0.18em,topsep=0.35em,leftmargin=2em}

\definecolor{deep}{RGB}{42,58,73}
\definecolor{soft}{RGB}{245,247,249}
\definecolor{line}{RGB}{145,154,163}
\definecolor{accent}{RGB}{104,76,55}
\definecolor{greenish}{RGB}{57,92,76}
\definecolor{warn}{RGB}{136,68,63}
\definecolor{purple}{RGB}{87,72,116}

% ── 双主题图解语义色彩（LaTeX 打印高对比白底）──
\definecolor{themebg}{HTML}{FFFFFF}
\definecolor{themefg}{HTML}{111111}
\definecolor{themegray}{HTML}{666666}
\definecolor{themecurve}{HTML}{1D63B8}
\definecolor{themealert}{HTML}{C53030}
\definecolor{themeamber}{HTML}{B86E00}
\definecolor{themeblue}{HTML}{1D63B8}

\newenvironment{handbookdiagram}{\par\vspace{0.4em}\noindent\begin{center}}{\end{center}\vspace{0.4em}}
\newcommand{\diagramcaption}[1]{\par\vspace{0.3em}{\small\textbf{#1}}}

\pgfdeclarelayer{background}
\pgfdeclarelayer{foreground}
\pgfsetlayers{background,main,foreground}

\hypersetup{
  colorlinks=true,
 linkcolor=deep,
  urlcolor=accent,
  pdftitle={大数定律与中心极限定理}
}

\pagestyle{fancy}
\fancyhf{}
\lhead{\small\color{deep}\textbf{数学一 · 概率统计 Topic 06}}
\rhead{\small\color{deep}依概率收敛 · 大数定律 · 中心极限定理 · 渐近分布}
\cfoot{\small\thepage}
\renewcommand{\headrulewidth}{0.4pt}

\newtcolorbox{corebox}[1]{breakable,colback=soft,colframe=deep,title=\textbf{#1},boxrule=0.8pt,arc=2mm}
\newtcolorbox{methodbox}[1]{breakable,colback=white,colframe=greenish,title=\textbf{#1},boxrule=0.7pt,arc=1.5mm}
\newtcolorbox{warnbox}[1]{breakable,colback=white,colframe=warn,title=\textbf{#1},boxrule=0.7pt,arc=1.5mm}
\newtcolorbox{examplebox}[1]{breakable,colback=white,colframe=purple,title=\textbf{#1},boxrule=0.7pt,arc=1.5mm}
\newcommand{\kw}[1]{\textbf{\color{deep}#1}}

\title{\vspace{-1.1cm}\textbf{大数定律与中心极限定理}\\[0.4em]
\large 随机性为什么能在大量重复中呈现确定性规律与正态渐近？}
\author{数学一 · 概率统计 Topic 06}
\date{}

\begin{document}
\maketitle

\begin{corebox}{母问题：随机性为什么能在大量重复中呈现稳定位置与可描述的剩余波动？}
先固定本册最小记号：i.i.d. 表示“相互独立且同分布”；$X_n\xrightarrow{P}X$ 表示对任意 $\varepsilon>0$ 都有 $P(\lvert X_n-X\rvert>\varepsilon)\to0$；$X_n\xrightarrow{d}X$ 表示分布函数在极限分布的连续点处收敛。下文把大数定律简称 LLN，把中心极限定理简称 CLT。后文的直觉解释都不得替代这些定义。

在经典的 i.i.d.、$E(X_i)=\mu$、$0<\operatorname{Var}(X_i)=\sigma^2<\infty$ 这一共同条件包下，同一组随机平均会出现两种\textbf{彼此相关但互不推出}的现象：
1. \kw{稳定位置：} $\bar X_n$ 靠近确定常数 $\mu$；大数定律回答“最终稳定到哪里”；
2. \kw{剩余波动：} 把误差按 $\sqrt n$ 的尺度放大后出现非退化极限；经典中心极限定理回答“稳定中心附近怎样波动”。

因此正确的结构不是“先 LLN，再由 LLN 推出 CLT”，而是从同一生成条件分别进入两条结论：
\[
\boxed{
\begin{array}{c}
X_1,\ldots,X_n\ \text{i.i.d.},\quad E(X_i)=\mu,\quad \operatorname{Var}(X_i)=\sigma^2\in(0,\infty)\\[0.5em]
\begin{array}{ccc}
\text{稳定问题} & & \text{波动问题}\\[-0.1em]
\bar X_n\xrightarrow{P}\mu & \qquad &
\dfrac{\sqrt n(\bar X_n-\mu)}{\sigma}\xrightarrow{d}N(0,1)\\
\text{大数定律} && \text{中心极限定理}
\end{array}
\end{array}}
\]
某些大数定律可以放宽有限方差条件，某些中心极限定理也允许非同分布；这些条件变化必须分别核对，不能把一条定理的资格借给另一条。
\end{corebox}

\section{本册定位与职责划分}

\begin{itemize}
  \item \kw{本册主干拥有（Owns）：}依概率/依分布等收敛语言、大数定律、中心极限定理、二项正态近似以及精确--渐近边界。
  \item \kw{拓展层：}几乎必然与 $L^p$ 收敛的更完整关系、Borel--Cantelli、Lyapunov/Lindeberg/Berry--Esseen、Slutsky、Delta 与 $o_P/O_P$ 用来解释机制和连接后续统计推断；它们不得伪装成数学一每道题都必须调用的主干清单。
  \item \kw{本册使用（Uses）：}Topic 05 的期望、方差与 Chebyshev 不等式；Topic 07 的样本统计量对象。
  \item \kw{本册不拥有（Does not own）：}小样本下精确抽样分布（Topic 07 负责）、统计量渐近估计性质（Topic 08 负责）。
  \item \kw{输出目标（Output）：}大样本下的概率近似计算与极限定理判断框架。
\end{itemize}

\section{第一层：收敛语言与强度边界}

设 $X_n,X$ 定义在同一概率空间上。四种常用收敛回答的问题不同：
\begin{itemize}
  \item 几乎必然收敛：$P(\{\omega:X_n(\omega)\to X(\omega)\})=1$，记作 $X_n\xrightarrow{a.s.}X$，是样本路径层面的结论；
  \item 依概率收敛：对所有 $\varepsilon>0$，$P(\lvert X_n-X\rvert>\varepsilon)\to0$，记作 $X_n\xrightarrow{P}X$，是偏离概率层面的结论；
  \item $L^p$ 收敛：$E\lvert X_n-X\rvert^p\to0$，记作 $X_n\xrightarrow{L^p}X$，其中 $p=2$ 是均方收敛；
  \item 依分布收敛：分布函数在极限分布的连续点收敛，记作 $X_n\xrightarrow{d}X$，不要求定义在同一概率空间。
\end{itemize}
基本蕴含链为
\[
X_n\xrightarrow{a.s.}X\Longrightarrow X_n\xrightarrow{P}X\Longrightarrow X_n\xrightarrow{d}X,
\qquad X_n\xrightarrow{L^p}X\Longrightarrow X_n\xrightarrow{P}X.
\]
若 $q>p>0$，$L^q$ 收敛推出 $L^p$ 收敛；若极限为常数 $c$，则 $X_n\xrightarrow{d}c$ 与 $X_n\xrightarrow{P}c$ 等价。反向一般不成立，不能把“分布越来越像”写成“同一条样本路径越来越近”。

若 $g$ 在 $X$ 的实际取值处几乎处处连续，连续映射定理给出
\[
X_n\xrightarrow{P}X\Rightarrow g(X_n)\xrightarrow{P}g(X),\qquad
X_n\xrightarrow{d}X\Rightarrow g(X_n)\xrightarrow{d}g(X).
\]
不连续函数必须另行检查极限落在不连续点的概率。

\begin{handbookdiagram}
\resizebox{0.92\linewidth}{!}{%
% ── V14 四大随机收敛拓扑结构（大空间呼吸感 2 行布局） ──
\begin{tikzpicture}[
  scale=1.0,
  >=Stealth,
  font=\small,
  color=themefg,
  text=themefg
]

\pgfdeclarelayer{background}
\pgfdeclarelayer{foreground}
\pgfsetlayers{background,main,foreground}

% ══════════════════════════════════════════════════════════════
% ── 上半部第 1 行：左（拓扑结构主干） + 右（三大反例机制） ──
% ══════════════════════════════════════════════════════════════

% ── 1. 左卡片：四大收敛蕴含拓扑主干 ──
\begin{scope}[xshift=0cm, yshift=5.0cm]
  % 卡片底色
  \begin{pgfonlayer}{background}
    \fill[fill=themefg!4!themebg, draw=themegray!30, line width=0.6pt, rounded corners=5pt]
      (-0.2, -0.2) rectangle (7.1, 5.2);
  \end{pgfonlayer}

  % 标题
  \begin{pgfonlayer}{foreground}
    \node[anchor=north west] at (0.0, 4.95) {
      \large \textbf{\color{themecurve}1. 四大收敛蕴含拓扑主干}
    };
  \end{pgfonlayer}

  % 拓扑有向图
  \begin{pgfonlayer}{background}
    % 连线箭头
    \draw[color=themecurve, line width=1.1pt, ->] (0.7, 3.4) -- (1.2, 2.6);
    \draw[color=themeamber, line width=1.1pt, ->] (2.1, 3.4) -- (1.6, 2.6);
    \draw[color=themealert, line width=1.1pt, ->] (1.4, 2.0) -- (1.4, 1.2);

    % 常数等价反向虚线
    \draw[color=themeamber, line width=0.9pt, dashed, ->] (1.9, 0.9) to[bend right=45] (1.9, 2.3);
  \end{pgfonlayer}

  \begin{pgfonlayer}{foreground}
    % 四大节点
    \node[color=themecurve, align=center] at (0.7, 3.7) {\textbf{几乎处处}\\[0.1em]\scriptsize a.s.};
    \node[color=themeamber, align=center] at (2.1, 3.7) {\textbf{均方收敛}\\[0.1em]\scriptsize $L^2$};
    \node[color=themeblue, align=center] at (1.4, 2.3) {\textbf{依概率收敛}\\[0.1em]\scriptsize $P$};
    \node[color=themealert, align=center] at (1.4, 0.9) {\textbf{依分布收敛}\\[0.1em]\scriptsize $d$};
    \node[color=themeamber] at (2.4, 1.6) {\tiny 常数 $c$};

    % 右侧文字说明
    \node[anchor=north west, text width=4.0cm] at (2.9, 4.3) {
      \small \textbf{单向强弱层级：}\\[0.25em]
      \scriptsize $\bullet$ a.s. 强收敛与 $L^2$ 均方互不推出；\\[0.15em]
      \scriptsize $\bullet$ 二者均可单向推出依概率收敛；\\[0.15em]
      \scriptsize $\bullet$ 依概率收敛单向推出依分布收敛。\\[0.7em]
      \small \textbf{常数等价特权：}\\[0.25em]
      \scriptsize 若极限为确定常数 $c$，则\\[0.15em]
      \color{themecurve}$X_n \xrightarrow{d} c \iff X_n \xrightarrow{P} c$。
    };
  \end{pgfonlayer}
\end{scope}

% ── 2. 右卡片：三大经典反例机制 ──
\begin{scope}[xshift=7.7cm, yshift=5.0cm]
  % 卡片底色
  \begin{pgfonlayer}{background}
    \fill[fill=themefg!4!themebg, draw=themegray!30, line width=0.6pt, rounded corners=5pt]
      (-0.2, -0.2) rectangle (7.1, 5.2);
  \end{pgfonlayer}

  % 标题
  \begin{pgfonlayer}{foreground}
    \node[anchor=north west] at (0.0, 4.95) {
      \large \textbf{\color{themeamber}2. 逆命题不成立：三大经典反例}
    };
  \end{pgfonlayer}

  \begin{pgfonlayer}{foreground}
    % 三大反例排印
    \node[anchor=north west, text width=6.7cm] at (0.1, 4.3) {
      \small $\bullet$ \textbf{\color{themecurve}$P \nRightarrow \text{a.s.}$ 滑动阶梯区间}：\\[0.15em]
      \scriptsize 移动区间指示函数 $I_{A_n}$，区间宽度 $\to 0$ 使测度趋零，但点态无限振荡无法收敛；\\[0.45em]
      \small $\bullet$ \textbf{\color{themeamber}$\text{a.s.} \nRightarrow L^2$ 极端逃逸尖峰}：\\[0.15em]
      \scriptsize $P(X_n = n^2) = \frac{1}{n^2}$，以概率 1 收敛到 0，但二阶矩 $E[X_n^2]=n^2 \to \infty$ 发散；\\[0.45em]
      \small $\bullet$ \textbf{\color{themealert}$d \nRightarrow P$ 对称相反同分布}：\\[0.15em]
      \scriptsize $X \sim N(0,1)$，令 $X_n = -X$；边缘恒同分布，但差值 $|X_n - X| = 2|X| \ne 0$。
    };
  \end{pgfonlayer}
\end{scope}

% ══════════════════════════════════════════════════════════════
% ── 下半部第 2 行：全宽卡片（连续映射 CMT 与 Slutsky 运算代数） ──
% ══════════════════════════════════════════════════════════════
\begin{scope}[xshift=0cm, yshift=0.3cm]
  % 卡片底色
  \begin{pgfonlayer}{background}
    \fill[fill=themefg!4!themebg, draw=themegray!30, line width=0.6pt, rounded corners=5pt]
      (-0.2, -0.2) rectangle (14.8, 4.3);
  \end{pgfonlayer}

  % 标题
  \begin{pgfonlayer}{foreground}
    \node[anchor=north west] at (0.0, 4.05) {
      \large \textbf{\color{themealert}3. 极限代数工具：连续映射（CMT）与~Slutsky~定理}
    };
  \end{pgfonlayer}

  \begin{pgfonlayer}{foreground}
    % 左栏：CMT 与 Slutsky 核心公式
    \node[anchor=north west, text width=7.2cm] at (0.1, 3.55) {
      \small $\bullet$ \textbf{\color{themecurve}连续映射定理 (CMT)}：\\[0.15em]
      \quad 若 $g$ 连续：$X_n \xrightarrow{\ast} X \implies g(X_n) \xrightarrow{\ast} g(X)$ \quad ($\ast \in \{a.s., L^r, P, d\}$)；\\[0.35em]
      \small $\bullet$ \textbf{\color{themeamber}Slutsky 代数定理（考场大杀器）}：\\[0.15em]
      \quad 若 $X_n \xrightarrow{d} X$ 且 $Y_n \xrightarrow{P} c$，则：\\[0.2em]
      \quad \scriptsize $X_n + Y_n \xrightarrow{d} X + c, \quad X_n Y_n \xrightarrow{d} cX, \quad X_n/Y_n \xrightarrow{d} X/c$。
    };

    % 中间分割线
    \draw[color=themegray!30, line width=0.6pt] (7.4, 0.3) -- (7.4, 3.5);

    % 右栏：考场判别清单与避坑准则
    \node[anchor=north west, text width=7.0cm] at (7.7, 3.55) {
      \small \textbf{考研判别清单与高频避坑：}\\[0.35em]
      \scriptsize $\bullet$ \textbf{判别依概率收敛}：\\[0.15em]
      \quad 若 $E[X_n]\to c$ 且 $\operatorname{Var}(X_n)\to 0$，则 $X_n \xrightarrow{P} c$；\\[0.3em]
      \scriptsize $\bullet$ \textbf{Slutsky 适用边界警示}：\\[0.15em]
      \quad 若 $Y_n \xrightarrow{d} Y$（非常数），则 $X_n+Y_n \nRightarrow X+Y$！\\[0.3em]
      \scriptsize $\bullet$ \textbf{Riesz 子列定理}：依概率收敛必存在几乎处处收敛子列。
    };
  \end{pgfonlayer}
\end{scope}

% ══════════════════════════════════════════════════════════════
% ── 底部全图全局总结 (无框轻量排印) ──
% ══════════════════════════════════════════════════════════════
\begin{pgfonlayer}{foreground}
  \node[anchor=north, text=themefg, text width=14.8cm, align=center] at (7.3, -0.2) {
    \footnotesize \textbf{收敛拓扑真理：} 样本轨道点态约束（a.s.）与均方能量积分（$L^2$）汇流于依概率测度收敛（$P$），最后投影为宏观累积分布（$d$）；Slutsky 运算严格要求次要项依概率收敛到常数。
  };
\end{pgfonlayer}

\end{tikzpicture}
%
}
\end{handbookdiagram}

\section{第二层：Borel--Cantelli 与路径事件}

事件“无穷多次发生”记为
\[
\{A_n\ \mathrm{i.o.}\}=\limsup A_n=\bigcap_{m\geq1}\bigcup_{n\geq m}A_n.
\]
第一 Borel--Cantelli 引理：若 $\sum_nP(A_n)<\infty$，则 $P(A_n\ \mathrm{i.o.})=0$，不要求独立。第二引理：若 $A_n$ 相互独立且 $\sum_nP(A_n)=\infty$，则 $P(A_n\ \mathrm{i.o.})=1$。它们把概率和转成路径层面的“只发生有限次/发生无穷次”，也是强收敛证明的重要工具。

\section{第三层：大数定律的完整分层}

\subsection{从方差衰减到弱大数律}

若 $X_i$ 两两独立、$\operatorname{Var}(X_i)\leq C$，则
\[
\operatorname{Var}\left(\frac1n\sum_{i=1}^nX_i\right)
=\frac1{n^2}\sum_{i=1}^n\operatorname{Var}(X_i)\leq\frac Cn.
\]
Chebyshev 立刻给出
\[
\frac1n\sum_{i=1}^n(X_i-EX_i)\xrightarrow{P}0.
\]
\begin{methodbox}{弱大数律的证明骨架}
这里的关键不是“样本量大”，而是方差按 $1/n$ 衰减：先由独立性消掉和式中的协方差项，再用 Chebyshev 把方差上界转换成偏离概率上界。若只知道 $E\lvert X_1\rvert<\infty$ 而没有有限方差，则不能沿这条 Chebyshev 路径，而要调用 Khinchin 弱大数律，并明确它的前提更换了。
\end{methodbox}

\begin{examplebox}{真题对象映射：2022 Q九 的样本二阶矩}
若 $X_1,\ldots,X_n$ 独立同分布，$E(X_1^k)=\mu_k$（$k=1,2,3,4$），题目考察
\[
T_n=\frac1n\sum_{i=1}^nX_i^2
\quad\text{偏离}\quad \mu_2.
\]
先把 $W_i=X_i^2$ 看成新的 i.i.d. 变量，则 $E(W_i)=\mu_2$、$\operatorname{Var}(W_i)=\mu_4-\mu_2^2$。于是
\[
\operatorname{Var}(T_n)=\frac{\mu_4-\mu_2^2}{n},
\qquad
P(\lvert T_n-\mu_2\rvert\ge\varepsilon)\le
\frac{\mu_4-\mu_2^2}{n\varepsilon^2}.
\]
不能看到“样本均值”就把中心写成 $\mu_1$，也不能把该上界当成精确概率；核心动作是先识别极限对象是 $X^2$ 的总体均值。
\end{examplebox}
若 $X_i$ 独立同分布且 $E\lvert X_1\rvert<\infty$，Khinchin 弱大数律给出
\[
\bar X_n\xrightarrow{P}E(X_1).
\]
它只需要一阶绝对矩，不需要有限方差。若再有 i.i.d. 与有限一阶绝对矩，Kolmogorov 强大数律给出
\[
\boxed{\bar X_n\xrightarrow{a.s.}E(X_1)}.
\]
Bernoulli 频率定律只是强大数律的示性变量特例：独立重复事件的频率几乎必然趋于其概率。

\subsection{大数律的统计推论}

若 $E\lvert X\rvert^k<\infty$，则样本 $k$ 阶原点矩
\[
A_{k,n}=\frac1n\sum_{i=1}^nX_i^k\xrightarrow{a.s.}E(X^k).
\]
若 $E(X^2)<\infty$，则
\[
\frac1n\sum X_i^2\xrightarrow{a.s.}E(X^2),\qquad \bar X_n^2\xrightarrow{a.s.}[E(X)]^2,
\]
所以 $B_{2,n}=n^{-1}\sum(X_i-\bar X_n)^2\xrightarrow{a.s.}\sigma^2$，且 $S_n^2=nB_{2,n}/(n-1)$ 也相合。对固定 $x$，经验分布函数
\[
F_n(x)=\frac1n\sum_{i=1}^nI_{\{X_i\le x\}}
\]
几乎必然趋于 $F(x)$；更强的 Glivenko--Cantelli 结论为 $\sup_x\lvert F_n(x)-F(x)\rvert\xrightarrow{a.s.}0$。样本对象的 Owner 见 Topic 07，本册只提供极限机制。

\section{第四层：中心极限定理的机制与扩展}

\subsection{经典 i.i.d. CLT}

若 $X_i$ i.i.d.、$E(X_i)=\mu$、$0<\sigma^2<\infty$，则
\[
\boxed{\frac{S_n-n\mu}{\sigma\sqrt n}=\frac{\sqrt n(\bar X_n-\mu)}\sigma\xrightarrow{d}N(0,1)}.
\]
因此 $\bar X_n$ 的误差尺度是 $n^{-1/2}$；这是“平方根规律”的来源。LLN 研究 $\bar X_n$ 是否靠近 $\mu$，CLT 研究把剩余误差放大 $\sqrt n$ 后呈现何种非退化形状。

\begin{handbookdiagram}
\resizebox{0.92\linewidth}{!}{%
% ── V13 大数定律与中心极限双尺度（大空间呼吸感 2 行布局） ──
\begin{tikzpicture}[
  scale=1.0,
  >=Stealth,
  font=\small,
  color=themefg,
  text=themefg
]

\pgfdeclarelayer{background}
\pgfdeclarelayer{foreground}
\pgfsetlayers{background,main,foreground}

% ══════════════════════════════════════════════════════════════
% ── 上半部第 1 行：左（大数定律） + 右（中心极限定理） ──
% ══════════════════════════════════════════════════════════════

% ── 1. 左卡片：大数定律（LLN） ──
\begin{scope}[xshift=0cm, yshift=5.0cm]
  % 卡片底色
  \begin{pgfonlayer}{background}
    \fill[fill=themefg!4!themebg, draw=themegray!30, line width=0.6pt, rounded corners=5pt]
      (-0.2, -0.2) rectangle (7.1, 5.2);
  \end{pgfonlayer}

  % 标题
  \begin{pgfonlayer}{foreground}
    \node[anchor=north west] at (0.0, 4.95) {
      \large \textbf{\color{themecurve}1. 大数定律（LLN）：原尺度塌缩}
    };
  \end{pgfonlayer}

  % 图像与文字左右分栏
  \begin{pgfonlayer}{background}
    % 坐标轴
    \draw[color=themegray, line width=0.6pt, ->] (0.2, 1.2) -- (2.9, 1.2) node[anchor=west] {\tiny $\bar x$};
    \draw[color=themegray, line width=0.6pt, ->] (1.5, 0.9) -- (1.5, 4.1);

    % 曲线演进（n=1, n=10, n->oo）
    \draw[color=themegray!60, line width=0.8pt, domain=-1.2:1.2, samples=30]
      plot ({1.5 + \x}, {1.2 + 0.6*exp(-\x*\x/0.8)});
    \draw[color=themeamber, line width=1.0pt, domain=-0.9:0.9, samples=30]
      plot ({1.5 + \x}, {1.2 + 1.4*exp(-\x*\x/0.25)});
    \draw[color=themecurve, line width=1.5pt, domain=-0.45:0.45, samples=30]
      plot ({1.5 + \x}, {1.2 + 2.6*exp(-\x*\x/0.04)});

    \fill[color=themecurve] (1.5, 1.2) circle (2.2pt);
  \end{pgfonlayer}

  \begin{pgfonlayer}{foreground}
    \node[anchor=north, color=themecurve] at (1.5, 1.1) {\small $\mu$};
    \node[anchor=east, color=themegray] at (0.6, 1.7) {\tiny $n=1$};
    \node[anchor=east, color=themeamber] at (0.8, 2.5) {\tiny $n=10$};
    \node[anchor=west, color=themecurve] at (1.75, 3.6) {\tiny $n \to \infty$};

    % 右侧文字说明
    \node[anchor=north west, text width=3.8cm] at (3.1, 4.3) {
      \small \textbf{物理机制（退化塌缩）：}\\[0.25em]
      \scriptsize $\operatorname{Var}(\bar X_n) = \frac{\sigma^2}{n} \to 0$；\\[0.15em]
      宏观随机涨落完全消失，概率质量集中塌缩于常数单点 $\mu$。\\[0.7em]
      \small \textbf{收敛命题：}\\[0.25em]
      \scriptsize 弱定律：$\bar X_n \xrightarrow{P} \mu$\\[0.15em]
      强大数：$\bar X_n \xrightarrow{a.s.} \mu$
    };
  \end{pgfonlayer}
\end{scope}

% ── 2. 右卡片：中心极限定理（CLT） ──
\begin{scope}[xshift=7.7cm, yshift=5.0cm]
  % 卡片底色
  \begin{pgfonlayer}{background}
    \fill[fill=themefg!4!themebg, draw=themegray!30, line width=0.6pt, rounded corners=5pt]
      (-0.2, -0.2) rectangle (7.1, 5.2);
  \end{pgfonlayer}

  % 标题
  \begin{pgfonlayer}{foreground}
    \node[anchor=north west] at (0.0, 4.95) {
      \large \textbf{\color{themeamber}2. 中心极限（CLT）：$\sqrt{n}$ 尺度展开}
    };
  \end{pgfonlayer}

  % 图像与文字左右分栏
  \begin{pgfonlayer}{background}
    % 阴影区间
    \fill[fill=themecurve!20!themebg, domain=-0.75:0.75]
      (0.75, 1.2) -- plot[domain=-0.75:0.75, samples=30] ({1.5 + \x}, {1.2 + 2.5*exp(-\x*\x/0.7)}) -- (2.25, 1.2) -- cycle;

    % 坐标轴
    \draw[color=themegray, line width=0.6pt, ->] (0.2, 1.2) -- (2.8, 1.2) node[anchor=west] {\tiny $z$};
    \draw[color=themegray, line width=0.6pt, ->] (1.5, 0.9) -- (1.5, 4.0);

    % 标准正态曲线
    \draw[color=themecurve, line width=1.4pt, domain=-1.25:1.25, samples=40]
      plot ({1.5 + \x}, {1.2 + 2.5*exp(-\x*\x/0.7)});

    % 分界虚线
    \draw[color=themeamber, line width=0.6pt, dashed] (0.75, 1.2) -- (0.75, 2.3);
    \draw[color=themeamber, line width=0.6pt, dashed] (2.25, 1.2) -- (2.25, 2.3);
  \end{pgfonlayer}

  \begin{pgfonlayer}{foreground}
    \node[anchor=north, color=themefg] at (1.5, 1.1) {\tiny $0$};
    \node[anchor=north, color=themeamber] at (0.75, 1.1) {\tiny $-1$};
    \node[anchor=north, color=themeamber] at (2.25, 1.1) {\tiny $1$};
    \node[anchor=south, color=themecurve] at (1.5, 3.75) {\scriptsize $N(0,1)$};

    % 右侧文字说明
    \node[anchor=north west, text width=3.9cm] at (3.0, 4.3) {
      \small \textbf{物理机制（显微展开）：}\\[0.25em]
      \scriptsize 乘以 $\frac{\sqrt{n}}{\sigma}$ 放大残差波动；\\[0.15em]
      微观残余涨落展开为普适对称的高斯钟形曲线。\\[0.7em]
      \small \textbf{收敛命题：}\\[0.25em]
      \scriptsize $Z_n = \frac{\sqrt{n}(\bar X_n-\mu)}{\sigma} \xrightarrow{d} N(0,1)$
    };
  \end{pgfonlayer}
\end{scope}

% ══════════════════════════════════════════════════════════════
% ── 下半部第 2 行：全宽卡片（双尺度核心对比与考场决策清单） ──
% ══════════════════════════════════════════════════════════════
\begin{scope}[xshift=0cm, yshift=0.3cm]
  % 卡片底色
  \begin{pgfonlayer}{background}
    \fill[fill=themefg!4!themebg, draw=themegray!30, line width=0.6pt, rounded corners=5pt]
      (-0.2, -0.2) rectangle (14.8, 4.2);
  \end{pgfonlayer}

  % 标题
  \begin{pgfonlayer}{foreground}
    \node[anchor=north west] at (0.0, 3.95) {
      \large \textbf{\color{themealert}3. 双尺度对比辨析与考场动作清单}
    };
  \end{pgfonlayer}

  \begin{pgfonlayer}{foreground}
    % 左栏：对比维度
    \node[anchor=north west, text width=7.2cm] at (0.1, 3.25) {
      \small $\bullet$ \textbf{\color{themecurve}研究对象}：$\bar X_n$（宏观）~vs~$Z_n$（显微）；\\[0.45em]
      \small $\bullet$ \textbf{\color{themeamber}极限形态}：退化常数~$\mu$~vs~普适高斯~$N(0,1)$；\\[0.45em]
      \small $\bullet$ \textbf{\color{themealert}核心功能}：证明参数估计一致性与宏观稳态。
    };

    % 中间分割线
    \draw[color=themegray!30, line width=0.6pt] (7.4, 0.3) -- (7.4, 3.4);

    % 右栏：收敛模式与考场警示
    \node[anchor=north west, text width=7.0cm] at (7.7, 3.25) {
      \small $\bullet$ \textbf{\color{themecurve}收敛模式}：依概率收敛~$\xrightarrow{P}$~vs~依分布收敛~$\xrightarrow{d}$；\\[0.45em]
      \small $\bullet$ \textbf{\color{themeamber}功能定位}：参数点估计~vs~渐近置信区间估计；\\[0.45em]
      \small $\bullet$ \textbf{\color{themealert}考场避坑警示}：\\[0.15em]
      \quad 求和区间概率必用~CLT~查表，勿用~LLN！
    };
  \end{pgfonlayer}
\end{scope}

% ══════════════════════════════════════════════════════════════
% ── 底部全图全局总结 (无框轻量排印) ──
% ══════════════════════════════════════════════════════════════
\begin{pgfonlayer}{foreground}
  \node[anchor=north, text=themefg, text width=14.8cm, align=center] at (7.3, -0.2) {
    \footnotesize \textbf{双尺度全景哲学：} 大数定律回答 \textbf{“收敛到何处”}（原尺度退化为确定常数 $\mu$）；中心极限定理回答 \textbf{“残余波动如何分布”}（$\sqrt{n}$ 尺度展开为高斯分布）；二者互不替代。
  };
\end{pgfonlayer}

\end{tikzpicture}
%
}
\end{handbookdiagram}

\begin{methodbox}{CLT 的前提映射}
题目中的样本均值必须先映射为 $S_n-n\mu$：中心化去掉确定位置，独立同分布与 $0<\sigma^2<\infty$ 提供经典 CLT 的适用条件，除以真实尺度 $\sigma\sqrt n$ 才得到标准正态极限。若方差未知，先由 Topic 05/06 证明 $S_n\xrightarrow{P}\sigma$，再用 Slutsky；这只产生渐近结论，不能把一般总体的 Student 化统计量改写成有限样本 $t$ 分布。
\end{methodbox}

\subsection{二项近似与连续性修正}

若 $S_n\sim\operatorname{Bin}(n,p)$，则
\[
\frac{S_n-np}{\sqrt{np(1-p)}}\xrightarrow{d}N(0,1).
\]
若用连续正态分布近似离散二项概率，连续性修正通常把整数边界移动 $0.5$，以减少离散—连续边界误差；例如
\[
P(S_n\le k)\approx
\Phi\!\left(\frac{k+0.5-np}{\sqrt{np(1-p)}}\right).
\]
这是近似路径的精度增强，不是 CLT 本身的额外定理条件。若题目明确要求“利用正态近似”或需要数值精度，应优先修正；若只考标准化结构、选项本身不区分 $0.5$ 修正，必须说明采用的考试约定。$np$ 与 $n(1-p)$ 不能过小；规则阈值只是经验提示，不是定理条件，偏斜时应报告近似风险。

\begin{examplebox}{精确正态抽样与 CLT 近似不是同一条路}
1998 Q十四 的总体本身为 $N(3.4,6^2)$，因此无论 $n$ 多大都有有限样本精确结论
\[
\bar X_n\sim N\left(3.4,\frac{36}{n}\right),
\qquad
P(1.4<\bar X_n<5.4)=2\Phi\left(\frac{\sqrt n}{3}\right)-1.
\]
要求概率至少 $0.95$，用 $\Phi^{-1}(0.975)=1.96$ 得 $\sqrt n/3\ge1.96$，故 $n\ge35$。

2020 Q八 则是 $S\sim\operatorname{Bin}(100,1/2)$ 的离散计数，题目明确要求 CLT 近似：未经修正的结构标准化为 $(55-50)/5=1$；若追求离散到连续的数值精度，再把边界改为 $55.5$ 得 $1.1$。前者是精确正态总体的有限样本分布，后者是一般离散计数的渐近近似，触发条件不能互换。
\end{examplebox}

\begin{examplebox}{真题精度回检：2020 Q八 的三种输出不能混写}
令 $S\sim\operatorname{Bin}(100,1/2)$，题目目标为 $P(S\le55)$。精确二项模型直接给
\[
P(S\le55)=\sum_{k=0}^{55}\binom{100}{k}2^{-100}\approx0.8643735.
\]
若按题目要求使用 CLT，$ES=50$、$\operatorname{SD}(S)=5$，不作连续性修正时
\[
P(S\le55)\approx\Phi\left(\frac{55-50}{5}\right)=\Phi(1)\approx0.8413447.
\]
把离散事件 $S\le55$ 的格点边界移到连续区间端点 $55.5$，得到连续性修正
\[
P(S\le55)\approx\Phi\left(\frac{55.5-50}{5}\right)=\Phi(1.1)\approx0.8643339.
\]
三者分别是\kw{精确值、渐近结构值、精度增强的渐近值}；不能把其中任一个写成另外两个的等号。实际答题先服从题目指定的输出层，再用数值回检判断近似方向和精度。
\end{examplebox}

\begin{methodbox}{Poisson 近似：稀疏计数的另一条路径}
若 $S_n\sim\operatorname{Bin}(n,p_n)$，$p_n\to0$ 且 $np_n\to\lambda$，固定 $k$ 时
\[
P(S_n=k)=\binom nkp_n^k(1-p_n)^{n-k}
\longrightarrow e^{-\lambda}\frac{\lambda^k}{k!}.
\]
所以当题目给出“大次数、单次概率小、$np$ 适中”并明确要求 Poisson 近似时，先取 $\lambda=np$，再在 Poisson 模型中求和；例如 2025 Q九中 $n=20,p=0.1$，$\lambda=2$，故
\[
P(S_n\le1)\approx e^{-2}\left(1+2\right)=3e^{-2}.
\]
这与 CLT 正态近似的识别对象不同：Poisson 路径保留稀疏计数的离散结构，正态路径则中心化并按 $\sqrt{np(1-p)}$ 放大。不能仅因样本量大就在二者之间任意替换；先读题目指定的近似机制，再做参数和尾事件回检。
\end{methodbox}

\subsection{非同分布与误差界}

对独立非同分布变量，令 $s_n^2=\sum_i\operatorname{Var}(X_i)$。Lyapunov 条件
\[
\frac1{s_n^{2+\delta}}\sum_iE\lvert X_i-EX_i\rvert^{2+\delta}\to0
\]
推出标准化和依分布收敛到 $N(0,1)$；Lindeberg 条件则直接排除少数极端项主导总方差。i.i.d. 且三阶绝对中心矩有限时，Berry--Esseen 给出
\[
\sup_x\lvert P((S_n-n\mu)/(\sigma\sqrt n)\le x)-\Phi(x)\rvert=O(n^{-1/2}).
\]
这些是机制与误差边界，不应伪装成所有题目的必需假设。

\section{第五层：Slutsky、Delta 与渐近阶}

若 $X_n\xrightarrow{d}X$、$Y_n\xrightarrow{P}c$（常数），则
\[
X_n+Y_n\xrightarrow{d}X+c,\qquad X_nY_n\xrightarrow{d}cX,
\qquad X_n/Y_n\xrightarrow{d}X/c\ (c\ne0).
\]
由大数律 $S_n\xrightarrow{P}\sigma$，经典 CLT 与 Slutsky 联立得到一般总体的 Student 化渐近结论：
\[
\frac{\sqrt n(\bar X_n-\mu)}{S_n}\xrightarrow{d}N(0,1).
\]
这不是一般总体的有限样本 $t$ 分布；精确 $t$ 归 Topic 07。

若 $\sqrt n(T_n-\theta)\xrightarrow{d}N(0,\tau^2)$，且 $g$ 在 $\theta$ 处可微，则一阶 Delta 方法为
\[
\boxed{\sqrt n[g(T_n)-g(\theta)]\xrightarrow{d}N(0,[g'(\theta)]^2\tau^2)}.
\]
它来自 $g(T_n)=g(\theta)+g'(\theta)(T_n-\theta)+o_P(n^{-1/2})$；若 $g'(\theta)=0$，必须改用二阶展开。

渐近记号
\[
X_n=o_P(a_n)\iff X_n/a_n\xrightarrow{P}0,
\]
而 $X_n=O_P(a_n)$ 表示 $X_n/a_n$ 在概率意义下有界。它们用来判断 Taylor 余项能否忽略，不能替代具体的收敛方式。

\section{第六层：精确、渐近与失效边界}

\begin{center}
\begin{tabularx}{\linewidth}{L{4.2cm}Y}
\toprule
场景 & 结论\\
\midrule
正态总体、$\sigma$ 已知 & $\sqrt n(\bar X-\mu)/\sigma$ 对任意 $n$ 精确为 $N(0,1)$\\
正态总体、$\sigma$ 未知 & $\sqrt n(\bar X-\mu)/S$ 对任意 $n$ 精确为 $t(n-1)$，由 Topic 07 负责抽样分布\\
一般总体、有限方差 & 标准化均值渐近为 $N(0,1)$，替换 $S$ 后仍需说明是渐近\\
重尾或无限方差 & 经典 CLT 可能失效，稳定分布或其他尺度可能出现\\
有依赖样本 & 需额外的遍历、混合、鞅或稳健条件，不能直接套 i.i.d. 定理\\
\bottomrule
\end{tabularx}
\end{center}

典型边界是 Cauchy 样本：均值不存在，样本均值对每个 $n$ 仍为 Cauchy，不会因样本量增大而稳定到有限常数。即使 $X_n\xrightarrow{d}X$，也不能自动推出 $E(X_n)\to E(X)$，还需要一致可积性或支配条件。

\section{第七层：统一选择流程}

\begin{enumerate}
  \item 先辨认目标：趋于常数是 LLN/相合性，标准化误差分布是 CLT，函数变换是连续映射或 Delta。
  \item 核对结构：独立、同分布、矩存在、方差是否正、样本是否有依赖。
  \item 写出中心化量和真实尺度：$S_n-n\mu$、$\sqrt n(\bar X_n-\mu)$ 或 $s_n^{-1}\sum(X_i-EX_i)$。
  \item 区分精确分布与渐近分布；“样本量大”只说明可能近似，不自动制造正态性。
  \item 最后校验收敛类型、极限对象和近似误差，明确哪些结论交给 Topic 07/08。
\end{enumerate}

\begin{warnbox}{高频边界}
依概率收敛一般不能推出几乎必然收敛；依分布收敛到非退化变量一般不能推出依概率收敛；有限均值不能保证经典 CLT；用 $S$ 替换 $\sigma$ 后只有正态总体才有有限样本精确 $t$；Delta 方法的一阶导数为零时必须换尺度。
\end{warnbox}

\section{贯穿母例：Bernoulli 频率的稳定与剩余波动}

\begin{examplebox}{母例：成功率的两种尺度}
设 $X_i\overset{\mathrm{i.i.d.}}\sim\operatorname{Bern}(p)$，$0<p<1$，令 $\bar X_n=n^{-1}\sum_iX_i$。
由 $E(X_i)=p$、$\operatorname{Var}(X_i)=p(1-p)$，
\[
\operatorname{Var}(\bar X_n)=\frac{p(1-p)}n,
\qquad \bar X_n\xrightarrow{P}p.
\]
这回答 LLN 的问题：频率最终稳定到 $p$。若把剩余误差放大，
\[
\frac{\sqrt n(\bar X_n-p)}{\sqrt{p(1-p)}}\xrightarrow{d}N(0,1),
\]
这回答 CLT 的问题：稳定位置附近的波动如何分布。取 $g(x)=x(1-x)$，当 $p\in(0,1)$ 时 Delta 方法还给出
\[
\sqrt n\{g(\bar X_n)-g(p)\}\xrightarrow{d}
N\!\left(0,(1-2p)^2p(1-p)\right).
\]
若 $p=0$ 或 $1$，方差退化，不能直接套用这条标准化公式；若 $p=1/2$，则 $g'(p)=0$，一阶 Delta 只给退化极限，若要观察非退化波动必须改用二阶展开。这就是先检查参数边界和导数阶数的必要性。
\end{examplebox}

\section{一页压缩}

\begin{corebox}{一句话总结}
\[
\boxed{\text{大数定律保证样本均值稳定收敛到期望，中心极限定理保证标准化放大后的剩余波动服从正态。}}
\]
\end{corebox}

\end{document}
